% =====================================================================
%  src/preamble.tex  —  导言区（Lead 独占；队友只读，禁止修改）
%
%  本文件定义的数学宏（队友只能用这些，禁止自造）：
%    \dd   -> 微分 d            \ee   -> 自然常数 e         \ii  -> 虚数 i
%    \Tr \SL \GL \End \Aut \id \Hom \Sym \ad \rank
%    \R \C \Z \Q \N \PP \HH \AA  -> 黑板粗体
%    \cA \cB \cD \cE \cF \cG \cH \cL \cM \cN \cO \cP \cS \cU \cV \cW \cX -> 花体
%    \ol{#1} -> 上划线          \eps -> \varepsilon
%    \figomit{编号}{图注}       -> 图从略标记
% =====================================================================
\usepackage[a4paper,margin=2.6cm,top=3.0cm,bottom=2.8cm]{geometry}
\usepackage{amsmath,amssymb,amsthm,mathtools}
\usepackage{mathrsfs}
\usepackage{graphicx}
\usepackage{array,multirow,booktabs,longtable}
\usepackage{caption}
\usepackage{enumitem}
\usepackage{xcolor}
\usepackage{fancyhdr}

% ---------- 字体（Windows，显式路径） ----------
\setCJKmainfont[Path=C:/Windows/Fonts/,BoldFont=simhei.ttf,ItalicFont=simkai.ttf]{simsun.ttc}
\setCJKsansfont[Path=C:/Windows/Fonts/]{simhei.ttf}
\setCJKmonofont[Path=C:/Windows/Fonts/]{simfang.ttf}
\setCJKfamilyfont{zhkai}[Path=C:/Windows/Fonts/]{simkai.ttf}
\setCJKfamilyfont{zhhei}[Path=C:/Windows/Fonts/]{simhei.ttf}
\setCJKfamilyfont{zhfang}[Path=C:/Windows/Fonts/]{simfang.ttf}
\providecommand{\kaishu}{\CJKfamily{zhkai}}
\providecommand{\heiti}{\CJKfamily{zhhei}}
\providecommand{\fangsong}{\CJKfamily{zhfang}}

% ---------- 数学宏 ----------
\newcommand{\dd}{\mathrm{d}}
\newcommand{\ee}{\mathrm{e}}
\newcommand{\ii}{\mathrm{i}}
\DeclareMathOperator{\Tr}{Tr}
\DeclareMathOperator{\SL}{SL}
\DeclareMathOperator{\GL}{GL}
\DeclareMathOperator{\End}{End}
\DeclareMathOperator{\Aut}{Aut}
\DeclareMathOperator{\id}{id}
\DeclareMathOperator{\Hom}{Hom}
\DeclareMathOperator{\Sym}{Sym}
\DeclareMathOperator{\ad}{ad}
\DeclareMathOperator{\rank}{rank}
\DeclareMathOperator{\Res}{Res}
\newcommand{\R}{\mathbb{R}}
\newcommand{\C}{\mathbb{C}}
\newcommand{\Z}{\mathbb{Z}}
\newcommand{\Q}{\mathbb{Q}}
\newcommand{\N}{\mathbb{N}}
\newcommand{\PP}{\mathbb{P}}
\newcommand{\HH}{\mathbb{H}}
% \AA 是 LaTeX 预定义命令（Å），本书要的是黑板粗体 A，故重定义
\renewcommand{\AA}{\mathbb{A}}
\newcommand{\cA}{\mathcal{A}}
\newcommand{\cB}{\mathcal{B}}
\newcommand{\cC}{\mathcal{C}}
\newcommand{\cD}{\mathcal{D}}
\newcommand{\cE}{\mathcal{E}}
\newcommand{\cF}{\mathcal{F}}
\newcommand{\cG}{\mathcal{G}}
\newcommand{\cH}{\mathcal{H}}
\newcommand{\cL}{\mathcal{L}}
\newcommand{\cM}{\mathcal{M}}
\newcommand{\cN}{\mathcal{N}}
\newcommand{\cO}{\mathcal{O}}
\newcommand{\cP}{\mathcal{P}}
\newcommand{\cS}{\mathcal{S}}
\newcommand{\cU}{\mathcal{U}}
\newcommand{\cV}{\mathcal{V}}
\newcommand{\cW}{\mathcal{W}}
\newcommand{\cX}{\mathcal{X}}
\newcommand{\eps}{\varepsilon}
\newcommand{\ol}[1]{\overline{#1}}
\newcommand{\figomit}[2]{\par\smallskip\noindent\begingroup\small\kaishu\color{gray}〔原文图 #1 从略：#2〕\endgroup\par\smallskip}
% 无编号插图（原书未标图号）：\figomitn{图注}
\newcommand{\figomitn}[1]{\par\smallskip\noindent\begingroup\small\kaishu\color{gray}〔原文插图（无编号）从略：#1〕\endgroup\par\smallskip}

% ---------- 定理类环境（编号一律手工给） ----------
\theoremstyle{definition}
\newtheorem*{definition}{定义}
\theoremstyle{plain}
\newtheorem*{theorem}{定理}
\newtheorem*{proposition}{命题}
\newtheorem*{lemma}{引理}
\newtheorem*{corollary}{推论}
\theoremstyle{definition}
\newtheorem*{example}{例}
\theoremstyle{remark}
\newtheorem*{remark}{注}

\makeatletter
\def\thmhead@plain#1#2#3{%
  \thmname{#1}\thmnumber{\@ifnotempty{#2}{ #2}}%
  \thmnote{\@ifnotempty{#3}{\hspace{0.35em}#3}}}
\let\thmhead\thmhead@plain
\makeatother

% ---------- 版式 ----------
\setlength{\parindent}{2em}
\linespread{1.18}
\setlist{nosep,leftmargin=2.2em}
\renewcommand{\arraystretch}{1.15}

\pagestyle{fancy}
\fancyhf{}
\fancyhead[L]{\small\kaishu 共形场论与规范对称性}
\fancyhead[R]{\small\thepage}
\renewcommand{\headrulewidth}{0.4pt}
\fancypagestyle{plain}{\fancyhf{}\fancyfoot[C]{\small\thepage}\renewcommand{\headrulewidth}{0pt}}

\ctexset{
  section = {
    format = \Large\bfseries\heiti,
    beforeskip = 2.2ex plus 1ex minus .2ex,
    afterskip  = 1.4ex plus .2ex,
    name = {第,章},
    number = \arabic{section},
  },
  subsection = {
    format = \large\bfseries\heiti,
    beforeskip = 1.8ex plus .6ex minus .2ex,
    afterskip  = 1.0ex plus .2ex,
  },
}

\usepackage[hidelinks,bookmarksnumbered=false,bookmarksopen=true,
            pdftitle={共形场论与规范对称性（中文译注本）},
            pdfauthor={Kenji Ueno（上野健爾）}]{hyperref}
\pdfstringdefDisableCommands{\def\ol#1{#1}}
